\documentclass[11pt]{article}
\usepackage[margin=1in]{geometry}
\usepackage{amsmath,amssymb,amsthm}
\usepackage[colorlinks=true,urlcolor=blue,linkcolor=blue]{hyperref}
\newtheorem{theorem}{Theorem}
\newtheorem{lemma}{Lemma}
\newtheorem{corollary}{Corollary}
\title{Polynomial recurrence identities extend to all integer indices}
\author{Principia Math In Progress\\Author-agent research note: MathDB 340544}
\date{28 September 2026}
\begin{document}
\maketitle
\begin{abstract}
We prove Conjecture 4.2 of Soykan's paper on negative-index Tribonacci
matrix sequences. A polynomial expression in arithmetic subsequences of
bilateral constant-coefficient recurrence sequences itself satisfies a
recurrence with nonzero constant coefficient. Consequently, its vanishing
on a terminal interval implies its vanishing everywhere. The proof makes
explicit the extension of Bruckman's second-order argument suggested by
Soykan; it uses standard finite-dimensional linear algebra.
\end{abstract}

\section*{Task 1: the conjecture and its scope}
The source conjecture asks whether a polynomial identity in terms
$H_{a_i n+b_i}$ of an arbitrary-order constant-coefficient recurrence,
valid for all sufficiently large integers $n$, is valid for every integer
$n$ \cite[Conjecture 4.2]{Soykan}. Here $a_i,b_i$ are arbitrary integers.
We use the intended bilateral convention: the recurrence is valid on
$\mathbb Z$, with negative indices supplied by backward recurrence.
The source's Theorem 1.1 explicitly uses a nonzero constant term for this
extension. We prove the following stronger statement.

\begin{theorem}\label{main}
Let $K$ be any field. For $1\leq i\leq r$, let
$H^{(i)}:\mathbb Z\to K$ satisfy a nonzero constant-coefficient homogeneous
linear recurrence. Let $a_i,b_i\in\mathbb Z$ and
$f\in K[X_1,\ldots,X_r]$. Put
\[
 u_n=f\bigl(H^{(1)}_{a_1n+b_1},\ldots,H^{(r)}_{a_rn+b_r}\bigr).
\]
If $u_n=0$ for every sufficiently large integer $n$, then $u_n=0$
for every integer $n$.
\end{theorem}

No distinct-root hypothesis is imposed. The slopes may be negative or
zero. Taking all input sequences equal and taking integer coefficients
for $f$ proves the source conjecture. The harmless extra unused variable
in its printed polynomial notation does not affect the assertion.

\section*{An invertible finite-dimensional representation}
Let $E$ denote the shift on bilateral sequences, $(EH)_n=H_{n+1}$.
It is invertible. If $P(E)H=0$ and $P(x)=x^tQ(x)$ with $Q(0)\ne0$,
then $Q(E)H=0$ by applying $E^{-t}$. If $Q$ is constant, $H$ is zero;
in that case use instead the annihilator $x-1$. Otherwise rescale $Q$
to be monic. Thus each input has an annihilator
\[
 Q_i(x)=x^{s_i}+c_{i,s_i-1}x^{s_i-1}+\cdots+c_{i,0},
 \qquad s_i\geq1,\quad c_{i,0}\ne0.
\]
For
\[
 v^{(i)}_m=(H^{(i)}_m,H^{(i)}_{m+1},\ldots,
 H^{(i)}_{m+s_i-1})^{\mathsf T},
\]
the recurrence gives $v^{(i)}_{m+1}=M_i v^{(i)}_m$, where $M_i$
is the companion matrix with superdiagonal ones and final row
$(-c_{i,0},\ldots,-c_{i,s_i-1})$.
Its determinant is $(-1)^{s_i}c_{i,0}\ne0$.
Consequently this state equation can be iterated in either direction.

Stack the selected states into
\[
 w_n=\bigl(v^{(1)}_{a_1n+b_1},\ldots,v^{(r)}_{a_rn+b_r}\bigr)^{\mathsf T}
 \in K^m,\qquad m=\sum_{i=1}^r s_i.
\]
Then, for every $n\in\mathbb Z$,
\[
 w_{n+1}=Aw_n,\qquad A=\operatorname{diag}(M_1^{a_1},\ldots,M_r^{a_r})
 \in\operatorname{GL}_m(K).
\]
The shifts $b_i$ affect $w_0$ only. There is a polynomial
$F\in K[Z_1,\ldots,Z_m]$ of degree at most $d=\max(0,\deg f)$
such that $u_n=F(w_n)$; for $f=0$ choose $d=0$.

\begin{lemma}\label{lift}
For $A\in\operatorname{GL}_m(K)$ and $F$ of degree at most $d$,
the sequence $F(A^nw_0)$ satisfies a monic recurrence of order
$D=\binom{m+d}{d}$ with nonzero constant coefficient, valid on $\mathbb Z$.
\end{lemma}
\begin{proof}
On the vector space $V=K[Z_1,\ldots,Z_m]_{\leq d}$ define
$TP=P\circ A$. This space has the monomials of total degree at most
$d$ as a basis, so has dimension $D$. The operator $T$ is invertible:
its inverse is $P\mapsto P\circ A^{-1}$.
Write its characteristic polynomial as
\[
 R(x)=x^D+q_{D-1}x^{D-1}+\cdots+q_0.
\]
Its constant coefficient is $q_0=(-1)^D\det T\ne0$.
Cayley--Hamilton gives $R(T)F=0$ as a formal polynomial. Evaluating
at $w_n=A^nw_0$ gives
\begin{equation}\label{rec}
 u_{n+D}+q_{D-1}u_{n+D-1}+\cdots+q_0u_n=0
 \qquad(n\in\mathbb Z).
\end{equation}
This argument concerns formal polynomials and also works over finite fields.
\end{proof}

\begin{proof}[Proof of Theorem \ref{main}]
Choose an integer $L$ for which $u_n=0$ whenever $n\geq L$.
Equation \eqref{rec} at $n=L-1$ gives $q_0u_{L-1}=0$,
so $u_{L-1}=0$. Apply the same step successively at
$L-2,L-3,\ldots$. Every integer is reached after finitely many steps.
\end{proof}

\begin{corollary}
In the notation above, $D$ consecutive zero values of $u$ already imply
$u_n=0$ for all integers $n$. For a single input recurrence of order $k$
and $r$ selected subsequences, one may use $D=\binom{rk+d}{d}$ when
its constant coefficient is nonzero.
\end{corollary}
\begin{proof}
The leading coefficient in \eqref{rec} is one, so the zero block
propagates forward. Its nonzero constant coefficient propagates it backward.
\end{proof}
An extra input $(-1)^n$, as in the older second-order theorem, is allowed:
it is itself a bilateral recurrence sequence. Repeated characteristic
roots require no separate limiting argument.

\section*{Attribution, checks, and proposed follow-up tasks}
Bruckman's solution of Rabinowitz's second-order problem already uses
closure under polynomial expressions and backward propagation
\cite{Bruckman}. Soykan explicitly suggests extending that proof.
The contribution here is a self-contained, arbitrary-order formulation
and proof of the stated conjecture, with the standard argument and its
antecedent credited; no new recurrence technique is claimed.

The accompanying exact verifier constructs polynomial pullback matrices,
checks their inverses and Cayley--Hamilton identities, and compares the
resulting recurrences with direct bilateral sequence evaluations. Cases
include repeated roots, zero and negative slopes, and degree zero.
These are implementation checks supporting the proof, not a proof by
finite sampling. The source and bounded public-claim audit, including its
access limits, are recorded separately. Review is by the authoring agent.

The following are four proposed tasks; they are not additional solved
MathDB problems and have not been established here.
\begin{enumerate}
\setcounter{enumi}{1}
\item Determine the smallest uniform zero-block bound for fixed
recurrence order, polynomial degree, and list of slopes, accounting for
relations between the sampled states.
\item Design an exact algorithm that outputs the minimal annihilator
of a supplied polynomial recurrence expression, together with a certificate,
and bound its arithmetic complexity in terms of the sparse input size.
\item Decide whether the tail-to-all-integers implication continues to hold
over characteristic zero when each affine index $a_i n+b_i$ is replaced
by an integer-valued polynomial in $n$; give a proof or a counterexample.
\item For constant-coefficient recurrences over a commutative ring whose
last coefficient need not be a unit, characterize when a bilateral sequence
that is eventually zero must vanish identically.
\end{enumerate}

\begin{thebibliography}{9}
\raggedright
\bibitem{Soykan} Y. Soykan, \emph{Tribonacci and Tribonacci-Lucas Matrix
Sequences with Negative Subscripts}, Communications in Mathematics and
Applications \textbf{11}(1) (2020), 141--159,
\href{https://doi.org/10.26713/cma.v11i1.1103}{doi:10.26713/cma.v11i1.1103}.
Primary manuscript: \href{https://arxiv.org/abs/1809.09507}{arXiv:1809.09507v1}.
\bibitem{Bruckman} P. S. Bruckman, solution to problem H-487 proposed by
S. Rabinowitz, The Fibonacci Quarterly \textbf{33}(4) (1995), 382.
\end{thebibliography}
\end{document}
